\section{Introduction}
\label{sec:intro}

Image generation has advanced rapidly in recent years, evolving from
specialized text-to-image models~\cite{sd3,flux} to unified multimodal
models (UMMs)~\cite{bagel,chameleon,januspro,emu3,emu3.5,illume+,metaquery}
that integrate visual understanding and generation within a single
architecture. Reinforcement learning has emerged in parallel as an
effective post-training stage~\cite{flowgrpo,dancegrpo,awm,diffusionnft},
consistently lifting compositional fidelity and instruction-following. The success of an RL recipe, however, rests on two
complementary pillars. 
The optimization algorithm governs the stability of long-horizon training,
while the reward model determines the ceiling that the trained policy can
ultimately approach. 
However, designing a practical reward model that remains both efficient and reliable is still challenging.


Recent studies have made substantial progress in reward modeling for image generation.
One line of work builds reward models from large-scale
human preference annotations, using these data to align visual-text
representations or vision-language models with human judgments~\cite{pickscore,imagereward,hpsv2,rewarddance,unifiedreward}. While effective, these methods depend on expensive annotation, difficult data collection, and costly iteration. Another line of work bypasses preference training by reusing pretrained MLLMs as zero-shot reward sources~\cite{viescore,li2025uniworld,alphagrpo}. Direct scalar or logit-based feedback is training-free, but can be sensitive to judge calibration and scoring noise~\cite{viescore,li2025uniworld}. Question-decomposition pipelines improve reliability by verifying atomic prompt requirements one by one, but introduce substantial engineering complexity~\cite{alphagrpo}. This leaves open the design of an efficient, preference-label-free, and off-the-shelf reward model for open-sourced image generators.


We propose \method, a training-free reward function that turns any
pretrained MLLM into an image-generation reward model by reusing its
image-text alignment knowledge, without preference labels or additional training. 
Given a generated image, \method conditions a frozen MLLM on the image and performs a single teacher-forced forward pass on the prompt, 
producing a token-level likelihood profile that reflects how well each textual requirement can be read from the image. 
We refer to this profile as the \emph{semantic spectrum} of the image-text pair, 
and aggregate it into a scalar reward using the mean prompt-token log-likelihood. 
In this way, \method uses the prompt-likelihood spectrum induced by the model's own image-conditioned language modeling objective as a deterministic, off-the-shelf reward signal, as shown in~\autoref{fig:teaser}~(a), instead of asking the MLLM to judge an image or answer decomposed questions.

We further specialize \method within unified multimodal models (UMMs),
yielding a special case we call \emph{Self-\method}. Since a UMM integrates image
understanding and image generation into a single framework, its 
understanding branch could directly provide the reward signals for its
generation branch.
In this way, the model improves its own generation
using its own understanding ability, without relying on any external
model or external knowledge, forming a closed-loop self-improving framework in which the policy evolves by
leveraging its own intrinsic capability. 
Beyond removing external dependencies, this design has a structural advantage. 
The understanding branch shares the same tokenizer, vision encoder, and pretraining
distribution as the generation branch, so the way it reads an image is naturally aligned with 
the distribution from which the generation branch samples images. 
We therefore hypothesize that this intrinsic alignment makes the policy's own understanding branch a
particularly suitable reward source for its own generation.

\begin{figure*}[t]
  \centering
  \includegraphics[width=\linewidth]{figures/teaser.pdf}
    % \vspace{-5pt}
  \caption{\textbf{Overview of \method.}
  (a) Pretrained MLLMs naturally induce a semantic spectrum that measures how well a generated image aligns with the prompt.
  \method aggregates this into a reward for T2I RL.
  (b) During RL training, \method steadily increases together with visible improvements in complex scene generation.
  (c) We study nine reward MLLM backbones from four MLLM families, with external reward MLLMs spanning three families and 4B to 235B parameters. 
  Scaling the reward MLLM backbones brings non-monotonic gains.
  Qwen3-VL-30B-A3B achieves the best performance among external MLLMs,
  while Self-\method, using BAGEL's own understanding branch as the reward model, outperforms all external MLLMs.
  (d) Both \method and Self-\method bring significant and consistent improvements across all six downstream benchmarks compared to the baselines.
  }
  \label{fig:teaser}
  \vspace{-10pt}
\end{figure*}

We validate \method through a broad image-generation RL study covering
two diffusion models, three RL training algorithms, nine reward MLLM backbones from four MLLM
families ranging from 4B to 235B parameters, and evaluate real generalization improvements on five out-of-distribution downstream text-to-image generation benchmarks.
On both SD3.5-M~\cite{sd3} and BAGEL~\cite{bagel}, \method brings significant improvements on all downstream benchmarks, showing the generalization ability of the proposed prompt-likelihood reward.
With BAGEL as the policy model, both \method and Self-\method consistently improve the BAGEL baseline by +10.0 on TIIF-Bench~\cite{tiif} and +4.3/5.5 on GenEval~\cite{geneval}, respectively.
\method also outperforms the strong RL baseline, AlphaGRPO~\cite{alphagrpo}, by 6.3 and 2.1 gains on TIIF-Bench and GenEval, respectively.
Moreover, Self-\method exhibits better improvements compared to \method with similar-scale MLLM backbones, and comparable or even better performance compared to 4x or 30x larger MLLMs,
e.g., +1.2 on GenEval, +2.1 on GenEval2~\cite{kamath2025geneval}, and +2 on WISE~\cite{wise} compared to \method with the best MLLM backbone.
These results support our hypothesis that distributional alignment between the reward model and the policy can be
as important as the scale of the reward model. The contributions of the paper are summarized as follows:
% \begin{itemize}[nolistsep, leftmargin=*]
\begin{itemize}[leftmargin=*, itemsep=1pt, topsep=1pt]
  \item We propose \method, a training-free and off-the-shelf reward
  function that turns any pretrained MLLM into a reward model for text-to-image generation. 
  By calculating the image-conditioned prompt-token likelihood, 
  \method reuses pretrained image-text alignment without preference labels, 
  reward-model fine-tuning, scalar judging, or question decomposition.

  \item We introduce {Self-\method}, the special case in which the policy's own understanding branch in a unified multimodal model serves as the reward model. 
  This forms a closed-loop self-improving framework without external reward models or 
  external knowledge, where the model leverages its own understanding ability to improve generation performance.

  \item We conduct a broad image-generation RL study covering two
  generator backbones, three RL algorithms, nine reward MLLM backbones
  from four MLLM families spanning $4$B--$235$B parameters, and five
  downstream benchmarks. The results show the effectiveness of \method and Self-\method, and reveal that the benefits of scaling the reward model are not monotonic
  and that the well-aligned self-reward model can match or even outperform much larger
  external models.
\end{itemize}
